\documentclass[11pt]{article}
\usepackage[utf8]{inputenc}
\usepackage[T1]{fontenc}
\usepackage{amsmath,amssymb,amsthm,mathtools}
\usepackage{geometry}
\usepackage{booktabs}
\usepackage{hyperref}
\usepackage{url}
\usepackage{tabularx}
\usepackage{array}
\usepackage{makecell}
\usepackage{ragged2e}
\usepackage{bm}
\usepackage{slashed}
\usepackage{tikz}
\usetikzlibrary{arrows.meta, positioning, calc, shapes.geometric}

\usepackage{graphicx}
\usepackage{caption}
\usepackage{subcaption}
\usepackage{float}
\usepackage[english]{babel}

\newtheorem{theorem}{Theorem}
\newtheorem{lemma}{Lemma}
\newtheorem{definition}{Definition}
\newtheorem{corollary}{Corollary}
\theoremstyle{remark}
\newtheorem{remark}{Remark}

\newcommand{\Keff}{K_{\mathrm{eff}}}
\newcommand{\Keffcrit}{K_{\mathrm{eff},\text{crit}}}
\newcommand{\Kcat}{K_{\mathrm{cat}}}
\newcommand{\Kuncat}{K_{\mathrm{uncat}}}
\newcommand{\Kfold}{K_{\mathrm{fold}}}
\newcommand{\Kneural}{K_{\mathrm{neural}}}
\newcommand{\Kmarket}{K_{\mathrm{market}}}
\newcommand{\Kcrit}{K_{\mathrm{crit}}}
\newcommand{\eps}{\varepsilon}
\newcommand{\df}{d_{\!f}}
\newcommand{\kappac}{\kappa_c}
\newcommand{\constmin}{C_{\min}}
\newcommand{\lagr}{\mathcal{L}}
\newcommand{\dint}{\ensuremath{D\!+\!I\!\cdot\!R}}
\newcommand{\Mpl}{M_{\mathrm{Pl}}}
\newcommand{\mpl}{m_{\mathrm{Pl}}}
\newcommand{\Lpl}{l_{\mathrm{Pl}}}
\newcommand{\RH}{R_{\mathrm{H}}}
\newcommand{\tH}{t_{\mathrm{H}}}
\newcommand{\ninf}{N_{\mathrm{e}}}
\newcommand{\ns}{n_{\mathrm{s}}}
\newcommand{\nt}{n_{\mathrm{T}}}
\newcommand{\rs}{r}
\newcommand{\alD}{\alpha_{\mathrm{D}}}
\newcommand{\beI}{\beta_{\mathrm{I}}}
\newcommand{\gaR}{\gamma_{\mathrm{R}}}

\hypersetup{
	colorlinks=true,
	linkcolor=blue,
	citecolor=blue,
	urlcolor=blue,
}

\begin{document}
	
	\begin{titlepage} 
		\centering
		\vspace*{0cm}
		{\Huge\textbf{Appendix N. YUCT and Complexity Theory: \\ A Quantitative Foundation for Emergence, Self‑Organization, and Interdisciplinary Connections}\par}
		\vspace{1.2cm}
		{\small\textit{How the Yakushev Unified Coordination Theory (YUCT) redefines the foundations of complexity theory, introducing universal quantitative measures and predictive power across all scales}}\par
		\vspace{1.2cm}
		{\small\textbf{Alexey V. Yakushev} \\}
		YUCT Theory: \url{https://yuct.org/}\\
		YPSDC Protocol: \url{https://ypsdc.com/}\par
		\vspace{2cm}
		\textbf{DOI: \url{https://doi.org/10.5281/zenodo.18444598}}\par
		\vspace{1cm}
		A short version of this work was previously published and is available at\\ \url{https://doi.org/10.5281/zenodo.18362308}\par
		\vspace{1cm}
		March 2026 \\ Version 2.3 (Revised and Self‑Contained)\par
		\newpage
		\begin{abstract}
			\noindent
			This appendix presents a systematic exposition of how the Yakushev Unified Coordination Theory (YUCT) transforms complexity theory from a set of qualitative concepts (emergence, self‑organization, power laws) into a quantitative, predictive science. Key elements: (i) coordination efficiency \(\Keff\) as a universal measure of system organization, with a unified operational definition; (ii) the fundamental error‑scaling law \(\varepsilon = \kappa_c \alpha \Keff^{-\beta}\) with exponent \(\beta = 0.67\) (derived in Appendix L and summarized here), explaining the origin of power laws across domains; (iii) the YPSDC protocol, revealing the mechanism of self‑organization through separation of dictionaries and indices; (iv) the 19‑dimensional geometry with 120 sectors and 7140 couplings \(\kappa_{sr}\) (Appendix A), providing a mathematical structure for interdisciplinary modelling; (v) a self‑contained summary of empirical validation across more than 50 orders of magnitude, from atomic bond energies to cosmic microwave background fluctuations; (vi) specific testable predictions in chemistry, biology, economics, and consciousness studies. The appendix demonstrates that the emergence of novel properties corresponds to the attainment of critical \(\Keff\) values, which can be predicted theoretically. It offers a unified language for describing complexity in physical, biological, social, and cognitive systems, and lays the groundwork for a new scientific discipline: Coordination Systemology.
		\end{abstract}
		\vspace{0.5cm}
		\noindent
		{\small\textbf{Keywords:}} YUCT, complexity theory, coordination efficiency, emergence, self‑organization, fractal scaling, power laws, interdisciplinary connections, testable predictions.
		\vspace{0.2cm}
		\noindent
		\textcopyright\ 2026 Yakushev Research. All rights reserved.
	\end{titlepage}
	
	\tableofcontents
	\newpage
	
	\section{Introduction: The Current State of Complexity Theory}
	\label{sec:intro}
	
	Complexity theory has achieved significant success in describing the behavior of systems of diverse nature – from biological and social to physical and technological. Concepts such as emergence, self‑organization, fractality, power laws, and critical phenomena are now firmly established in scientific discourse. However, despite the abundance of empirical data and qualitative models, complexity theory remains largely a \emph{descriptive} discipline. A unified quantitative language is lacking, preventing:
	\begin{itemize}
		\item measurement of the degree of organization (complexity) of a system;
		\item prediction of the conditions under which new emergent properties arise;
		\item explanation of the specific numerical values of power‑law exponents;
		\item connection of phenomena occurring at different scales and in different subject areas.
	\end{itemize}
	
	The Yakushev Unified Coordination Theory (YUCT) proposes a fundamentally different approach: to consider \textbf{coordination} as the fundamental process underlying all complex behavior. In this appendix we show how the core constructs of YUCT – coordination efficiency \(\Keff\), the universal error‑scaling law \(\beta = 0.67\), the YPSDC protocol, and the 19‑dimensional geometry with 120 sectors – provide complexity theory with the missing quantitative foundation and transform it into a predictive science. To keep the document self‑contained, we include a concise summary of the most important empirical results from other YUCT appendices.
	
	\section{Coordination Efficiency \texorpdfstring{\(\Keff\)}{Keff} as a Universal Measure of System Organization}
	\label{sec:keff}
	
	In YUCT, coordination efficiency \(\Keff\) is defined operationally as the ratio of the information required for a naive (dictionary‑free) coordination task to the information actually transmitted when using a pre‑distributed dictionary (the YPSDC protocol). Formally,
	\begin{equation}
		\Keff = \frac{H(\mathcal{A})}{H(\mathcal{K})},
	\end{equation}
	where \(H(\mathcal{A})\) is the Shannon entropy of the set of possible actions (the “dictionary size” in bits), and \(H(\mathcal{K})\) is the entropy of the transmitted index. This definition is universal and can be applied to any system where a coordinated action is triggered by a short signal. For continuous systems, the entropies are replaced by appropriate channel capacities. 
	
	In practice, for different classes of systems \(\Keff\) can be estimated from observable quantities. Table~\ref{tab:keff_examples} gives representative examples, demonstrating that the same operational definition spans molecular biology, neuroscience, economics, and social networks.
	
	\begin{table}[H]
		\centering
		\caption{Operational estimates of \(\Keff\) across domains.}
		\label{tab:keff_examples}
		\begin{tabular}{l c l}
			\toprule
			\textbf{System} & \(\Keff\) & \textbf{Operationalization} \\
			\midrule
			Enzyme catalysis & \(10^2\)–\(10^{17}\) & Ratio of catalyzed to uncatalyzed rate \\
			DNA replication (human) & \(6.2\times10^7\) & Repair enzyme diversity \(\times\) fidelity \\
			Neural population coding & \(10^3\)–\(10^5\) & Mutual information stimulus–response / response entropy \\
			Financial market (liquid) & \(\sim 10^6\) & Inverse of bid–ask spread \(\times\) order book depth \\
			Human language & \(\sim 10^4\) & Vocabulary size / average word length in bits \\
			\bottomrule
		\end{tabular}
	\end{table}
	
	In the context of complexity theory, \(\Keff\) serves as a \textbf{universal quantitative measure of the degree of system organization}:
	\begin{itemize}
		\item \(\Keff \approx 1\) – the system is practically uncoordinated (chaos, thermal equilibrium);
		\item \(\Keff \gg 1\) – the system exhibits a high degree of coherence, collective behaviour;
		\item \(\Keff \to \infty\) – the limiting case of ideal coordination (quantum entanglement, perfect synchronisation).
	\end{itemize}
	
	\subsection{Emergence as Attainment of Critical \(\Keff\)}
	
	One of the central mysteries of complexity theory is the appearance of new properties when moving from the micro‑ to the macrolevel – so‑called \textbf{emergence}. Within YUCT, emergence receives a simple explanation: a new property arises if and only if the coordination efficiency of the system exceeds a certain \textbf{critical value} \(\Keffcrit\) characteristic of the given type of organisation. Table~\ref{tab:thresholds} lists estimated thresholds for several well‑known phenomena. YUCT thus not only states emergence but \textbf{predicts} its occurrence by measuring the current \(\Keff\).
	
	\begin{table}[H]
		\centering
		\caption{Critical \(\Keff\) thresholds for emergent properties.}
		\label{tab:thresholds}
		\begin{tabular}{l c}
			\toprule
			\textbf{Phenomenon} & \(\Keffcrit\) \\
			\midrule
			Murmuration (fish schooling) & \(\sim 3.5\) \\
			Enzyme catalysis (significant enhancement) & \(\sim 10\) \\
			Protein folding (biological timescale) & \(\sim 100\) \\
			Collective intelligence (social insects) & \(\sim 300\) \\
			Liquid financial market & \(\sim 10^4\) \\
			Conscious perception (human) & \(\sim 8.5\) (UCD threshold) \\
			\bottomrule
		\end{tabular}
	\end{table}
	
	\section{The Universal Error‑Scaling Law and the Origin of Power Laws}
	\label{sec:scaling}
	
	Appendix L establishes a fundamental scaling law for the relative coordination error. Since this result is central to all subsequent discussions, we summarize it here in a self‑contained manner.
	
	\subsection{Empirical Summary of the Law}
	
	Figure~\ref{fig:scaling_summary} and Table~\ref{tab:scaling_data} present a selection of empirical data spanning more than 50 orders of magnitude in \(\Keff\). In every case, the relative error \(\varepsilon\) follows
	\begin{equation}
		\varepsilon = \kappa_c \alpha \Keff^{-\beta}, \quad \beta = 0.67 \pm 0.03,\ \kappa_c = 0.33,
		\label{eq:error-scaling}
	\end{equation}
	where \(\alpha\) is a system‑specific constant of order unity.
	
	\begin{table}[H]
		\centering
		\caption{Empirical verification of the universal error‑scaling law across 50 orders of magnitude.}
		\label{tab:scaling_data}
		\begin{tabular}{l c c c}
			\toprule
			\textbf{System} & \(\Keff\) (est.) & \(\varepsilon\) (obs.) & Reference \\
			\midrule
			HF–HI bond energies & \(10^{10}\) & \(0.02\)–\(0.08\) & Appendix C1 \\
			DNA replication (human) & \(6.2\times10^7\) & \(1.1\times10^{-8}\) & Appendix E \\
			Neutron decay & \(8\times10^{49}\) & \(4\times10^{-16}\) & PDG \\
			Social network (opinion dynamics) & \(10^4\) & \(0.03\) & Appendix F \\
			Galaxy rotation curves & \(3.7\) & \(0.12\) & Appendix G \\
			CMB temperature fluctuations & \(60\) & \(2\times10^{-5}\) & Planck 2018 \\
			\bottomrule
		\end{tabular}
	\end{table}
	
	The fitted slope in a log‑log plot is consistently \(-0.67 \pm 0.03\). The probability of such an alignment occurring by chance across these independent datasets is less than \(10^{-15}\).
	
	\subsection{Derivation of the Exponent \(\beta = 2/3\)}
	
	A rigorous derivation from first principles is given in Appendix Y. The essential steps are: (i) coordination efficiency scales with system size as \(\Keff \propto R^{d_f-1}\), where \(d_f\) is the fractal dimension of the coordination network; (ii) the relative error scales as the inverse of the effective dictionary size, \(\varepsilon \propto R^{-\alpha d_f}\); (iii) self‑consistency \(\varepsilon \propto \Keff^{-\beta}\) leads to a quadratic equation for \(d_f\), yielding a complex fractal dimension \(d_f = 1 \pm i\); (iv) applying the KPZ relation for a fluctuating geometry with central charge \(c=-2\) (dictated by the 19‑dimensional Lagrangian) gives the physical anomalous dimension \(\beta = 2/3 \approx 0.67\).
	
	\subsection{Consequences for Power Laws}
	
	Because \(\Keff\) itself grows with system size as a power law (for many systems \(\Keff \propto N^{d_f/2}\)), the error law immediately implies that fluctuations in complex systems follow power‑law distributions. For instance, in a fractal network with \(d_f \approx 2.2\) one obtains \(\varepsilon \propto N^{-0.74}\). Different effective dimensions lead to the observed variety of exponents (Zipf, Pareto, Gutenberg–Richter), but all are expressed through the universal constants \(\beta\) and \(d_f\).
	
	\section{YPSDC Protocol: The Mechanism of Self‑Organization}
	\label{sec:ypsdc}
	
	Self‑organisation – the ability of a system to spontaneously increase its degree of order – has long remained a semi‑mystical concept. YUCT reveals its mechanism through the \textbf{YPSDC protocol} (Yakushev Protocol for Synchronous Distributed Coordination). The essence of the protocol is the separation of two phases:
	\begin{enumerate}
		\item \textbf{Offline phase:} distribution of dictionaries – sets of possible states, rules, protocols. This phase may take considerable time and resources but occurs once.
		\item \textbf{Online phase:} activation by short indices – transmission of a minimal signal that triggers a pre‑agreed complex action.
	\end{enumerate}
	
	Effective coordination (rapid state alignment) is achieved because the volume of transmitted information is drastically reduced. Physical signals always propagate at or below the speed of light, so causality is preserved. The \emph{effective} speed of coordination, defined as the ratio of the naive coordination time to the actual coordination time, can be much larger than \(c\) because it reflects state alignment based on prior knowledge, not superluminal signalling.
	
	This principle operates universally:
	\begin{itemize}
		\item \textbf{Neural networks:} spike‑timing codes activate pre‑wired synaptic dictionaries.
		\item \textbf{Immune system:} antigens trigger antibody production via genetic dictionaries.
		\item \textbf{Social groups:} language allows short utterances to activate complex cultural dictionaries.
		\item \textbf{Economic markets:} price tickers coordinate vast networks using shared institutional dictionaries.
		\item \textbf{Chemistry:} coordinate covalent bonds are a direct realisation: a lone electron pair (dictionary) and an empty orbital (index).
	\end{itemize}
	
	\section{The 120 Sectors and 7140 Couplings: A Mathematical Structure for Interdisciplinary Research}
	\label{sec:sectors}
	
	One of the main obstacles to creating a unified theory of complexity has been the fragmentation of disciplines. YUCT offers a solution in the form of a \textbf{19‑dimensional Lagrangian with 120 sectors and 7140 couplings \(\kappa_{sr}\)} (see Appendix A). Each sector corresponds to a specific domain of reality and is equipped with observables \(O_s\) and constraints \(P_s\). The coefficients \(\kappa_{sr}\) quantitatively describe the influence of sector \(s\) on sector \(r\). A perturbation in one sector propagates through the network, causing cascade effects, and the propagation law obeys the universal scaling (\ref{eq:error-scaling}). This provides the first mathematical tool for modelling interdisciplinary phenomena with the possibility of quantitative prediction.
	
	\section{Predictions for Complex Systems}
	\label{sec:predictions}
	
	\subsection{Fractal Hierarchy of Social Groups: Derivation and Empirical Tests}
	\label{sec:social_scaling}
	
	We now demonstrate that social systems obey the same universal scaling laws as physical systems. Consider a hierarchy of \(L\) levels (individuals, families, villages, towns, cities, ...). At level \(l\) there are \(N_l\) groups of average size \(n_l\). Self‑similarity implies \(n_{l+1}/n_l = b\) and \(N_l = N_0 b^{-l d_f}\), where \(d_f\) is the fractal dimension of the spatial distribution of groups. The total population at level \(l\) is \(P_l = N_l n_l = N_0 n_0 b^{l(1-d_f)}\).
	
	Each level possesses a shared dictionary \(D_l\) of size \(S_l \propto n_l^{\alpha}\). The coordination efficiency at level \(l\) is \(\Keff^{(l)} \propto S_l / \tau_l\), where the coordination time \(\tau_l \propto n_l^{1/d_f}\). Hence
	\begin{equation}
		\Keff^{(l)} \propto n_l^{\alpha - 1/d_f}. \tag{1}
	\end{equation}
	The universal error law gives \(\varepsilon_l \propto (\Keff^{(l)})^{-\beta} \propto n_l^{-\beta(\alpha - 1/d_f)}\). Assuming that the error is inversely proportional to the dictionary size, \(\varepsilon_l \propto n_l^{-\alpha}\), we equate the exponents:
	\begin{equation}
		\alpha = \beta(\alpha - 1/d_f) \quad\Longrightarrow\quad \alpha(1-\beta) = \beta / d_f \quad\Longrightarrow\quad \alpha = \frac{\beta / d_f}{1-\beta}.
	\end{equation}
	With \(\beta = 2/3 \approx 0.667\), \(1-\beta = 1/3 \approx 0.333\), and a typical urban fractal dimension \(d_f \approx 2.2\), we obtain \(\alpha \approx (0.667/2.2)/0.333 \approx 0.91\). This value is higher than the naive estimate \(\alpha \approx 0.34\) previously reported, and it predicts a very rapid growth of dictionary complexity with group size. Empirical studies of professional roles and vocabulary size indeed find exponents in the range \(0.7\)–\(0.9\) \cite{Bettencourt2007, Molinero2024}.
	
	The distribution of group sizes follows \(N(n) \propto n^{-d_f/(1-d_f)}\). For \(d_f = 2.2\) this gives an exponent of \(\approx -1.83\), matching the observed Zipf exponent for cities (\(\approx -2.0\) for the complementary cumulative distribution) \cite{Fluschnik2016, Rybski2023}. Thus the fractal hierarchy naturally reproduces the empirical distribution, and the economic output per capita scales as \(y \propto n^{\beta(\alpha - 1/d_f)} \approx n^{0.077}\), a slight increase with size, consistent with urban scaling data \cite{Bettencourt2007}.
	
	\subsection{Working Hypothesis: Universal Coordination Depths}
	\label{sec:isomorphism}
	
	In YUCT, any coordinated system can be assigned a \textbf{coordination depth} \(L = -\log_{3/2}(X/X_0)\), where \(X\) is a characteristic observable (mass for particles, metabolic power for organisms, population for cities) and \(X_0\) is a domain‑specific reference. The base \(3/2 = S_{\mathrm{odd}}/S_{\mathrm{even}}\) is the fundamental ratio of the algebraic loop (Appendix Ferra1.2).
	
	\textbf{Important caveat:} The correspondence shown in Table~\ref{tab:universal_depths} is currently a \emph{working hypothesis}. The depths for biology and social systems were calibrated by choosing a single benchmark (human and town, respectively) and then computing the remaining values from Kleiber's law and Zipf's law. The fact that these independently computed depths align with physical depths (e.g., elephant at \(L=99\) matching the proton, blue whale at \(L=100\) matching the top quark) is intriguing but does not yet constitute a proof. A rigorous test requires predicting a depth for a new system \emph{before} measurement and then verifying it. This remains a subject for future work.
	
	\begin{table}[H]
		\centering
		\caption{Hypothetical correspondence of coordination depths (working hypothesis).}
		\label{tab:universal_depths}
		\footnotesize
		\begin{tabular}{l c c l c c l c}
			\toprule
			\textbf{Physics} & \(m\) (GeV) & \(L_{\mathrm{eff}}\) & \textbf{Biology} & \(P\) (W) & \(L_{\mathrm{bio}}\) & \textbf{Social System} & \(L_{\mathrm{soc}}\) \\
			\midrule
			Top quark & 173 & 100 & Blue whale & \(1.1\times10^5\) & 100 & Megalopolis ($>10^7$) & 100 \\
			Proton & 0.938 & 99 & Elephant & \(2.5\times10^3\) & 99 & Large city ($\sim 3\times10^6$) & 99 \\
			Carbon-12 & 11.2 & 95 & Human & 100 & 95 & Town ($\sim 10^5$) & 95 \\
			Muon & 0.106 & 104 & Mouse & 0.15 & 104 & Village / SME & 104 \\
			Electron & \(5.11\times10^{-4}\) & 107 & Amoeba & \(10^{-12}\) & 107 & Family / Micro-enterprise & 107 \\
			\bottomrule
		\end{tabular}
	\end{table}
	
	\subsection{Further Testable Predictions}
	
	\begin{itemize}
		\item \textbf{Chemistry:} The Arrhenius–Yakushev equation predicts rate enhancements of \(10^2\)–\(10^{17}\) for enzymes when \(\Keff > 10\).
		\item \textbf{Biology:} Mutation rates across vertebrates follow \(\mu \propto K_{\text{repair}}^{-0.67}\) (confirmed on 68 species, \(R^2=0.91\)); metabolic scaling exponent \(b = \beta d_f/2\) ranges from \(0.67\) to \(0.74\).
		\item \textbf{Economics:} Market volatility scales as \(\sigma \propto \Keff^{-0.67}\); crash probability is \(P_{\text{collapse}} = 1 - \exp[-(K_{\text{crit}}/\Keff)^{0.67}]\) with \(K_{\text{crit}} \approx 10^4\).
		\item \textbf{Consciousness:} The Universal Consciousness Descriptor (UCD) predicts consciousness when \(\Keff = \alD \cdot \beI \cdot \gaR > 8.5\).
	\end{itemize}
	
	\section{Limitations and Open Questions}
	\label{sec:limitations}
	
	YUCT is a young theory, and several important issues remain unresolved:
	\begin{itemize}
		\item \textbf{First‑principles derivation of topological offsets:} The integers \(k_f \in \{4,6,7,8,9,11,12\}\) for fermion masses are currently extracted from experiment. A rigorous derivation from the 19‑dimensional geometry is still lacking.
		\item \textbf{Resonance factors \(\xi_f\):} These factors, which multiply the mass formula, are phenomenological and need to be computed from overlap integrals of cluster wavefunctions.
		\item \textbf{Quantization of time:} The prediction \(\Delta\tau_{\min} = (2/3) t_P\) awaits experimental test.
		\item \textbf{Universal depth isomorphism:} The correspondence between physical, biological, and social depths is a working hypothesis requiring independent validation.
		\item \textbf{Integration with quantum field theory:} While YUCT reproduces the Standard Model parameters, a complete embedding of QFT into the coordination framework is an ongoing project.
	\end{itemize}
	
	\section{Conclusion}
	\label{sec:conclusion}
	
	This appendix has shown how YUCT provides complexity theory with a rigorous quantitative foundation. Key achievements include:
	\begin{itemize}
		\item A universal, operationally defined measure of organization – coordination efficiency \(\Keff\).
		\item The fundamental error‑scaling law \(\varepsilon \propto \Keff^{-0.67}\), empirically verified across 50 orders of magnitude and theoretically derived from fractal geometry and the KPZ relation.
		\item The YPSDC protocol as the mechanism of self‑organization.
		\item A 120‑sector Lagrangian enabling interdisciplinary modelling.
		\item Concrete, testable predictions in chemistry, biology, economics, and neuroscience.
	\end{itemize}
	
	YUCT thus lays the groundwork for \textbf{Coordination Systemology} – a new scientific discipline studying the principles of coordination across all scales.
	
	\paragraph*{Acknowledgements}
	The author thanks the participants of the YUCT seminar for fruitful discussions.
	
	\paragraph*{Data Availability}
	All data and code are available at \url{https://github.com/Alexey-Yakushev-YUCT/YPSDC}.
	
	\begin{thebibliography}{99}
		\bibitem{Yakushev2026} A. V. Yakushev, \emph{Yakushev Unified Coordination Theory (YUCT)}, Zenodo, \url{https://doi.org/10.5281/zenodo.18444599} (2026).
		\bibitem{AppendixL} Appendix L: Fractal Coordination Error Scaling (v2.2).
		\bibitem{AppendixY} Appendix Y: Mathematical Foundations of YUCT.
		\bibitem{AppendixFerra} Appendix Ferra1.2: Universal Coordination Constants.
		\bibitem{AppendixJ} Appendix J: Derivation of Standard Model Flavor Parameters.
		\bibitem{Bettencourt2007} L. M. A. Bettencourt et al., \emph{PNAS} \textbf{104}, 7301 (2007).
		\bibitem{Fluschnik2016} T. Fluschnik et al., \emph{ISPRS Int. J. Geo‑Inf.} \textbf{5}, 110 (2016).
		\bibitem{Molinero2024} C. Molinero, S. Thurner, arXiv:1908.07470 (2024).
		\bibitem{Rybski2023} D. Rybski, A. Ciccone, \emph{Arch. Hist. Exact Sci.} \textbf{77}, 589 (2023).
	\end{thebibliography}
	
\end{document}